%Radio description
{\Large Радіодіапазони}
%Draw a sinewave with
%	{Initial Amplitude x.xxx}{Initial Time-width x.xxx}{Initial Number of cycles x}
%	{Second Amplitude x.xxx}{Second Time-width x.xxx}{Second Number of cycles x}
%	{Total repeats}
\newlength{\nAmplitude}		\newlength{\nAmplitudeNeg}
\newlength{\nTime}		\newlength{\nTimeTwo}		\newlength{\nTimeThree}		\newlength{\nTimeFour}
\newlength{\nSAmplitude}	\newlength{\nSAmplitudeNeg}
\newlength{\nSTime}		\newlength{\nSTimeTwo}		\newlength{\nSTimeThree}	\newlength{\nSTimeFour}
\newlength{\nTimeEnd}		\newlength{\nSTimeEnd}		\newlength{\nSTimeStart}
\newcommand{\sinewave}[7]{
	\setlength{\nAmplitude}{#1}
	\setlength{\nAmplitudeNeg}{\nAmplitude*-1}
	\setlength{\nTime}{#2}
	\setlength{\nTimeTwo}{\nTime*2}
	\setlength{\nTimeThree}{\nTime*3}
	\setlength{\nTimeFour}{\nTime*4}
	\setlength{\nTimeEnd}{\nTime*4*#3}
	\setlength{\nSAmplitude}{#4}
	\setlength{\nSAmplitudeNeg}{\nSAmplitude*-1}
	\setlength{\nSTime}{#5}
	\setlength{\nSTimeTwo}{\nSTime*2}
	\setlength{\nSTimeThree}{\nSTime*3}
	\setlength{\nSTimeFour}{\nSTime*4}
	\setlength{\nSTimeEnd}{\nTime*4*#3+\nSTime*4*#6}
	\multido{\dXTimePosition=0.000in+\nSTimeEnd}{#7}{
		\multido{\dXPosition=\dXTimePosition+\nTimeFour}{#3}{
			\rput(\dXPosition,0){
			\parabola(0.00,0)(\nTime,\nAmplitude)\parabola(\nTimeTwo,0)(\nTimeThree,\nAmplitudeNeg)
			}
		}
		\setlength{\nSTimeStart}{\dXTimePosition+\nTimeEnd}
		\multido{\dXPosition=\nSTimeStart+\nSTimeFour}{#6}{
			\rput(\dXPosition,0){
			\parabola(0.00,0)(\nSTime,\nSAmplitude)\parabola(\nSTimeTwo,0)(\nSTimeThree,\nSAmplitudeNeg)
			}
		}
	}
}
%

% None: f(x) = sin(20x)
% AM: f(x) = cos(.5x)∙sin(20∙x)
% FM: f(x) = cos(20∙x+15∙sin(x))


% Preamble:
% \pgfplotsset{width=7cm,compat=1.14}
% \begin{tikzpicture}
% \begin{axis}[
% hide x axis,
% hide y axis
% ]
% \addplot [red]{sin(deg(x))};
% \addplot [green]{cos(deg(x))};
% \end{axis}

% \begin{axis}[enlarge x limits=false]
% \addplot [red,samples=500] {sin(deg(x))};
% \addplot [orange,samples=7] {sin(deg(x))};
% \addplot [teal,const plot,samples=14]
% {sin(deg(x))};
% \end{axis}
% \end{tikzpicture}% <- eliminate space.
% \end{document}

% None: f(x) = sin(20x)
% AM: f(x) = cos(.5x)∙sin(20∙x)
% FM: f(x) = cos(20∙x+15∙sin(x))

\pgfplotsset{%
	axis line style={draw=none},
	tick style={draw=none},
	xticklabels=\empty,
	yticklabels=\empty,
	domain=0:3*pi,
	samples=1e3,
	width=2.3in,
	height=1in,
	xlabel shift=-8pt,
}

\begin{itemize}

%Country flags from here (using MIT license):
% https://github.com/stevenrskelton/flag-icon
% use inkscape to convert to eps which is smaller than using "convert"
%
%List of telecommunications regulatory bodies
% https://en.wikipedia.org/wiki/List_of_telecommunications_regulatory_bodies
%
\item {%
\setlength{\fboxsep}{0pt}%
\setlength{\fboxrule}{0.5pt}%
Радіоспектр (від ELF до EHF) містить набагато більше елементів, ніж можна показати на цій діаграмі. Тут наведено лише невелику вибірку діапазонів, що використовуються по всьому світу. Кожна країна має власні правила та норми розподілу діапазонів у цій області. FCC \fbox{\includegraphics[width=0.15in]{tex/us.eps}}, CRTC \fbox{\includegraphics[width=0.15in]{tex/ca.eps}}, OFCOM \fbox{\includegraphics[width=0.15in]{tex/gb.eps}}, BNA  \fbox{\includegraphics[width=0.15in]{tex/de.eps}}, CII  \fbox{\includegraphics[width=0.15in]{tex/cn.eps}}...
}%
\item Зв'язок за допомогою EMR здійснюється шляхом модуляції несучої хвилі:

\begin{tabular}{ccc}%
\begin{tikzpicture} % No modulation
	\begin{axis}[smooth,grid=none,scale=1,xlabel={Без модуляції},line width=1pt,]
		\addplot[white, smooth,xlabel={Без модуляції},line join=round, line cap=round] {sin(x*20*180/pi)};
	\end{axis}
\end{tikzpicture}&
\begin{tikzpicture} %AM
	\begin{axis}[smooth,grid=none,scale=1,xlabel={Амплітудна модуляція (AM)},line width=1pt,]
		%\addplot[white, smooth,line join=round, line cap=round] {sin(x*20*180/pi)*sin(x*1*180/pi)};
		\addplot[white, smooth,line join=round, line cap=round] {0.5*sin(x*20*180/pi)*(1+sin(x*1.7*180/pi))+8};

	\end{axis}
\end{tikzpicture}&
\begin{tikzpicture} % FM
	\begin{axis}[smooth,grid=none,scale=1,xlabel={Частотна модуляція (FM)},line width=1pt,]
		\addplot[white, smooth,line join=round, line cap=round] {sin((x*20 + 6*(sin(x*2*180/pi)))*180/pi)};
	\end{axis}
\end{tikzpicture}
\end{tabular}
%
% \item Not all references agree on the ULF band range, the HAARP range is used here.
%
\item {\bfseries РА}діо{\bfseries Д}етектування {\bfseries Т}а {\bfseries В}изначення {\bfseries В}ідстані (RADAR) використовує мікрохвилі для визначення відстані та швидкості об'єктів.

\item Радіо {\bfseries г}ромадського {\bfseries д}іапазону (CB) містить 40 станцій між 26.965 - 27.405MHz %\hspace{0.02in}
{% CB Radio
  \definecolor{BoxColor}{rgb}{0.6,0.68,0.86}
  \psframebox[linestyle=none,framesep=1pt,fillcolor=BoxColor,linearc=0]{\black CB}
}.
%
 Резонанси Шумана виникають у порожнині між Землею та іоносферою \hspace{0.02in}%
 \psframebox[fillstyle=none,linestyle=none]{\psscalebox{0.7}{\blip{0,-.07}{S}}}%
 \hspace{0.02in}.
%
 Водень випромінює радіодіапазонне EMR на довжині хвилі 21cm  \hspace{0.05in}%
 \psframebox[fillstyle=none,linestyle=none]{\psscalebox{0.7}{\blip{0,-.07}{H}}}%
 \hspace{0.02in}.
%
 Стандарти часу / частоти показані як \psframebox[linestyle=none] {\rput(0.1in,0){\timestandard}} \hspace{0.1in}.
%
 Інші окремі частоти представлені значками:\vspace{0.1in}\\
\begin{tabular}{cp{1.4in}cp{1in}cp{1.7in}}
%
\psframebox[linestyle=none]{\submarine{0.02,.0}{xxHz}}\hspace{0.2in}\vspace{0.05in} & Підводний човен&
\psframebox[linestyle=none]{\rput(0.15,.05)\pager}&Пейджер&
\psframebox[linestyle=none]{\rput(0.05,.05){\psframebox[fillstyle=solid,fillcolor=Itinerant,linecolor=Itinerant,linearc=0,framesep=1pt]{\textcolor{Black}{GMRS}}}}&Загальна мобільна радіослужба\\%
%
\psframebox[linestyle=none]{\rput(0.14,0.01){\psframebox[fillstyle=solid,fillcolor=green,linecolor=green,linearc=0]{\textcolor{Black}{xxm}}}}\hspace{.1in}\vspace{0.05in} & Аматорське / міжнародне&
\psframebox[linestyle=none]{\rput(0.14,.03){\weatherstation}\hspace{.03in}\vspace{0.08in}} & Погода&
\psframebox[linestyle=none]{\rput(0.15,.04){
	\psframe[linestyle=solid,linecolor=gray,fillstyle=solid,fillcolor=gray,linearc=0](-.2,-.08)(.2,.08)
	\psframe[hatchwidth=2pt, hatchsep=1.5pt,linestyle=solid,linecolor=yellow,fillstyle=hlines,hatchangle=45,hatchcolor=yellow,fillcolor=gray,linearc=0](-.2,-.08)(.2,.08)
	}
	\hspace{.2in}}\vspace{0.05in}
	& Короткі хвилі\\
%
%
\psframebox[linestyle=none]{\rput(0.11,.05){\psframebox[fillstyle=solid,fillcolor=Itinerant,linecolor=Itinerant,linearc=0,framesep=1pt]{\textcolor{Black}{FRS}}}}&Сімейна радіослужба&
\psframebox[linestyle=none]{\rput(0.15,-0.1){\wirelessmic}}&Бездрот. мік.&
\psframebox[linestyle=none]{
% 	\rput(0.15,.05){
% 		\psframebox[fillstyle=solid,linearc=0,linecolor=yellow,framesep=1pt,fillcolor=yellow,linewidth=1pt,linestyle=solid]{\textcolor{Black}{CP}}
% 	}
	\rput(0.08, 0.02){
		\psset{linestyle=solid, linewidth=0.02in}
		\psline[linecolor=cyan  ](-.15,0.04)(0.15,0.04)
		\psline[linecolor=green ](-.20,0.02)(0.10,0.02)
		\psline[linecolor=orange](-.10,0.00)(0.20,0.00)
	}
	\hspace{.03in}\vspace{0.08in}} & Мобільні телефони
	\psframebox[fillstyle=solid,linearc=0,linecolor=orange,framesep=1pt,fillcolor=orange,linewidth=1pt,linestyle=solid]{\textcolor{Black}{3G}}
	\psframebox[fillstyle=solid,linearc=0,linecolor=green, framesep=1pt,fillcolor=green, linewidth=1pt,linestyle=solid]{\textcolor{Black}{4G}}
	\psframebox[fillstyle=solid,linearc=0,linecolor=cyan,  framesep=1pt,fillcolor=cyan,  linewidth=1pt,linestyle=solid]{\textcolor{Black}{5G}}
\\
%
% \psframebox[linestyle=none]{
% 	\rput(0.11,.05){
% 		\psframebox[linestyle=none,framesep=1pt,fillcolor=red,linearc=0]{\white SOS}
% 	}
% }
% &
% 	\psset{dotsize=1pt 1}Distress signal, in Morse code:\vspace{0.18in}
% \rput(-1.4,-0.18){
% 	\rput(0,0.05){
% 		\psdots(0,0)(.1,0)(.2,0)\psline{cc-cc}(0.3,0)(0.42,0)\psline{cc-cc}(0.49,0)(0.61,0)\psline{cc-cc}(0.68,0)(0.8,0)\psdots(.9,0)(1,0)(1.1,0)
% 	}
% }%
\end{tabular}
\input{tex/morse.tex}
\end{itemize}
